\section{Основные определения}

\begin{definition}[Многокомпонентная контекстно-свободная грамматика]
    \emph{$m$-MCFG($r$)}~--- это четвёрка $\langle \Sigma, N, S, P \rangle$, где
    \begin{itemize}
        \item $\Sigma$~--- терминальный алфавит;
        \item $N$~--- конечное множество нетерминальных символов, каждому из которых приписан \emph{ранг} (арность, размерность) $k \ge 1$;
              максимальный ранг нетерминалов равен $m$;
        \item $S \in N$~--- стартовый нетерминал ранга $1$;
        \item $P$~--- конечное множество правил вида
              \[
              A(s_1,\ldots,s_k) \leftarrow B_1(x_1^1,\ldots,x_{k_1}^1), \ldots, B_n(x_1^n,\ldots,x_{k_n}^n),
              \]
              где
              \begin{itemize}
                  \item $A$~--- нетерминал ранга $k$, $B_i$~--- нетерминалы ранга $k_i$, $n \leq r$;
                  \item все переменные $x^i_j$ попарно различны;
                  \item каждый $s_i$~--- строка над $\Sigma \cup X$, где $X = \bigcup_{i=1}^n\bigcup_{j=1}^{k_i} \{x^i_j\}$.
              \end{itemize}
    \end{itemize}
    Если $m$ и $r$ несущественны или ясны из контекста, мы будем говорить просто \emph{MCFG}.
\end{definition}

Интуитивно, правило $A(s_1,\ldots,s_k) \leftarrow B_1(\overline{x}_1),\ldots,B_n(\overline{x}_n)$ означает, что если из каждого нетерминала $B_i$ выводится кортеж строк $\overline{w}_i$, то из $A$ выводится кортеж, полученный подстановкой $\overline{w}_i$ вместо переменных $x^i_j$ в строки $s_1,\ldots,s_k$.

\begin{definition}[Порождаемый язык]
    Язык, порождаемый MCFG $G = \langle \Sigma, N, S, P \rangle$, определяется как
    \[
    L(G) = \{ \omega \in \Sigma^* \mid S(\omega) \text{ выводимо в } G \}.
    \]
    Класс языков, порождаемых $m$-MCFG($r$), обозначается $m$-MCFL($r$).
    Совокупность всех MCFL обозначается MCFL.
\end{definition}

\begin{remark}
    При $m=1$ класс $1$-MCFL совпадает с классом контекстно-свободных языков, а $1$-MCFG в ослабленной нормальной форме Хомского~--- это в точности КС-грамматики в ОНФХ.
\end{remark}

\subsection{Примеры}

Для начала покажем, что любой контекстно-свободный язык является $1$-MCFL.
Приведём MCFG для известных нам контекстно-свободных языков.

\begin{example}
    Язык вложенных скобок $\{a^n b^n \mid n \ge 0\}$:
    \begin{align}
        S(axb) &\leftarrow S(x) \nonumber \\
        S(\varepsilon) &\leftarrow \nonumber
    \end{align}
\end{example}

\begin{example}
    Язык Дика на одном типе скобок:
    \begin{align}
        S(ax_1bx_2) &\leftarrow S(x_1), S(x_2) \nonumber \\
        S(\varepsilon) &\leftarrow \nonumber
    \end{align}
\end{example}

Теперь рассмотрим язык, не являющийся контекстно-свободным:
\[
L = \{a^n b^m c^n d^m \mid n, m \ge 0\}.
\]
Этот язык порождается следующей $2$-MCFG($2$):
\begin{align*}
    S(x_1 x_2 x_3 x_4) &\leftarrow P(x_1,x_3), Q(x_2,x_4) \\
    P(ax_1, cx_2) &\leftarrow P(x_1,x_2) \\
    P(\varepsilon,\varepsilon) &\leftarrow \\
    Q(bx_1, dx_2) &\leftarrow Q(x_1,x_2) \\
    Q(\varepsilon,\varepsilon) &\leftarrow
\end{align*}
Здесь нетерминал $P$ порождает пары $(a^n, c^n)$, нетерминал $Q$~--- пары $(b^m, d^m)$, а стартовое правило собирает их в нужном порядке.

\subsection{Разновидности MCFG}

В зависимости от ограничений на использование переменных в правилах, выделяют следующие подклассы MCFG~\sidecite{seki1991multiple}:

\begin{itemize}
    \item \textbf{Неудаляющая (non-erasing)}: каждая переменная $x^i_j$, встречающаяся в правой части правила, хотя бы один раз используется в строках $s_1,\ldots,s_k$.
    \item \textbf{Непереставляющая (non-permuting)}: для каждого $i$ порядок переменных $x^i_1,\ldots,x^i_{k_i}$ сохраняется в строках $s_1,\ldots,s_k$ (т.е. если $j < l$, то $x^i_j$ встречается левее $x^i_l$).
    \item \textbf{Well-nested}: неудаляющая, непереставляющая, и для любых двух различных нетерминалов $B_i, B_{i'}$ их переменные не перекрещиваются:
          \[
          s_1\cdots s_k \notin (\Sigma \cup X)^*\, x^i_j\, (\Sigma \cup X)^*\, x^{i'}_{j'}\, (\Sigma \cup X)^*\, x^i_{j+1}\, (\Sigma \cup X)^*\, x^{i'}_{j'+1}\, (\Sigma \cup X)^*.
          \]
\end{itemize}

\begin{example}
    Следующее правило является well-nested (несмотря на то, что переменные $B$ не идут подряд, никакие два различных нетерминала не содержат запрещённого паттерна чередования вида $B_j \ldots C_{j'} \ldots B_{j+1} \ldots C_{j'+1}$):
    \[
    A(x_1, z_1 z_2, x_2, y_1 y_2 y_3, x_3) \leftarrow B(x_1,x_2,x_3), C(y_1,y_2,y_3), D(z_1,z_2).
    \]
    Следующее правило \emph{не} является well-nested (переменные $B$ и $C$ перекрещиваются: чередование $B_1$ ($x_1$) \ldots $C_1$ ($y_1$) \ldots $B_2$ ($x_2$) \ldots $C_2$ ($y_2$) образует запрещённый паттерн):
    \[
    A(z_1, x_1, y_1, x_2, z_2, y_2, x_3, y_3) \leftarrow B(x_1,x_2,x_3), C(y_1,y_2,y_3), D(z_1,z_2).
    \]
\end{example}

Помимо MCFG, существуют и более выразительные расширения:

\begin{itemize}
    \item \textbf{PMCFG (Parallel MCFG)}: допускают дублирование переменных, например $A(x, ax) \leftarrow B(x)$.
    \item \textbf{simpleLMG}: допускают многократное использование одной переменной, например $A(x,x) \leftarrow B(x), C(x)$.
\end{itemize}

Известно, что эти классы образуют следующую иерархию выразительной силы:
\[
MCFL \subsetneq PMCFL \subsetneq simpleLMG = P.
\]
В частности, язык $\{a^{2^n} \mid n \ge 0\}$ порождается PMCFG (правила $S(xx) \leftarrow S(x)$, $S(a) \leftarrow$), но не является MCFL.

\subsection{Языки MIX и $O_n$}

Важным примером, демонстрирующим выразительную силу MCFL, являются языки MIX и $O_n$.

\begin{itemize}
    \item $MIX = \{\omega \in \{a,b,c\}^* \mid |\omega|_a = |\omega|_b = |\omega|_c\}$~--- язык строк с равным количеством символов $a$, $b$ и $c$.
    \item $O_2 = \{\omega \in \{a,\overline{a},b,\overline{b}\}^* \mid |\omega|_a = |\omega|_{\overline{a}},\ |\omega|_b = |\omega|_{\overline{b}}\}$.
    \item $O_n = \{\omega \in \{a_1,\overline{a_1},\ldots,a_n,\overline{a_n}\}^* \mid |\omega|_{a_i} = |\omega|_{\overline{a_i}}\ \text{для всех } i\}$.
    \item $MIX_n = \{\omega \in \{a_1,\ldots,a_n\}^* \mid |\omega|_{a_1} = \cdots = |\omega|_{a_n}\}$.
\end{itemize}

Долгое время оставалось открытым вопрос, является ли MIX многокомпонентным контекстно-свободным языком.
Положительный ответ был дан Сальвати~\sidecite{salvati:inria-00564552}: MIX является $2$-MCFL.
Ключевая идея доказательства~--- рациональная эквивалентность MIX и $O_2$: если один из них является MCFL, то и другой тоже, поскольку $2$-MCFL образуют полный AFL.

Более того, язык $O_n$ является $n$-MCFL~\sidecite{GEBHARDT202241}.
Этот факт имеет прямое применение в статическом анализе программ, где $O_n$ может моделировать согласованное количество открывающих и закрывающих скобок для $n$ различных типов контекстов (например, для одновременного \emph{context-sensitive} и \emph{data-dependence} анализа).
Хотя точное описание ограничений для подобных задач анализа кода требует \textbf{шафла языков Дика} (Interleaved Dyck language),
который не является MCFL, различные приближающие его MCFL используются решения соответствующих задач~\sidecite{10.1145/3704854}.

Другим интересным языком является \textbf{многомерный язык Дика}. Вопрос о существовании $2$-MCFG для трёхмерного языка Дика $D^3$ остаётся открытым, хотя и предпринимались попытки построения соответствующих грамматик~\sidecite{10.1007/978-3-662-59620-3_5}.
